\chapter{System Implementation}
This chapter describes the implemented architecture at the v0.0.6 freeze.
The v0.0.7 role-argument and benchmark corrections are described in
Chapter~\ref{chap:v007}.
\label{chap:implementation}

\section{Architecture}
AnsibleGuard v0.0.6 is not a renamed ansible-lint configuration. It is an
orchestrator and semantic analyser with its own parsers, artifact model,
execution graph, rule registry, comparison engine, canonical data model, and
report generators. ansible-lint and GitLeaks are isolated evidence providers.

\begin{figure}[H]
\centering
\begin{tikzpicture}[
 c/.style={draw=ansblue,rounded corners=3pt,fill=anspaper,
 text width=32mm,minimum height=11mm,align=center,font=\small},
 a/.style={-{Latex[length=2.2mm]},thick,ansblue}]
\node[c] (entry) at (0,0) {CLI / Python API\\repository selection};
\node[c] (parse) at (4.3,0) {Discovery and tolerant\\YAML/Jinja parsing};
\node[c] (graph) at (8.6,0) {Artifact, symbol and\\execution graphs};
\node[c] (providers) at (0,-2.6) {Isolated ansible-lint\\and GitLeaks providers};
\node[c] (canon) at (4.3,-2.6) {Canonical occurrences\\and provider health};
\node[c] (rules) at (8.6,-2.6) {Internal semantic rules\\and source evidence};
\node[c] (outputs) at (4.3,-5.2) {Output-v3 validation\\JSON, CSV, SARIF, text};
\draw[a] (entry)--(parse); \draw[a] (parse)--(graph);
\draw[a] (graph)--(rules);
\draw[a] (entry)--(providers);
\draw[a] (rules)--(canon);
\draw[a] (providers)--(canon);
\draw[a] (canon)--(outputs);
\end{tikzpicture}
\caption{Implementation evidence flow. Internal rules and external providers
produce separate evidence that is reconciled in the canonical output model;
provider failures remain explicit health records.}
\label{fig:implementation_architecture}
\end{figure}

\section{Parsing and Artifact Model}
Discovery applies repository exclusions and records every selected path. YAML
is loaded tolerantly so Ansible tags such as \code{!vault} and \code{!unsafe}
do not abort the repository. Source coordinates are retained for findings.
Jinja expressions are analysed as a second language embedded in scalar values.
The classifier distinguishes playbooks, task and handler files, defaults,
role variables, inventory, group/host variables, metadata, and import-only
files. This prevents, for example, task-shape rules from running on inventory.

Jinja names introduced by \code{for} are lexical locals. A loop variable is
visible only between its matching \code{for} and \code{endfor}; nested blocks
use a scope stack. It neither leaks after the block nor becomes a repository
definition. Guards such as \code{x is defined}, short-circuit expressions,
\code{default}, and \code{omit} are attached to the expression scope in which
they protect a use.

\section{Repository Execution Semantics}
Static \code{include_tasks} and \code{import_tasks} targets form directed
execution edges. Persistent definitions, including \code{set_fact}, propagate
only along reachable static edges and only after their producer in program
order. They are never unioned across all repository files. This is important:
a fact created in one disconnected role must not suppress a finding in another.
Dynamic includes cannot safely name one successor and therefore contribute an
explicit uncertainty reason rather than a fabricated definition.

For a variable use $u$ in file $f$, let $D(u)$ be visible definitions, $G(u)$
lexically valid guards, $R(u)$ modelled runtime sources, and $Q(u)$ named
uncertainty sources. The state function is:
\[
S(u)=
\begin{cases}
\mathrm{CONFIRMED} & D(u)\cup G(u)\cup R(u)\neq\varnothing,\\
\mathrm{ORDERED\_CONFLICT} & \text{a producer exists but follows }u,\\
\mathrm{UNCERTAIN} & Q(u)\neq\varnothing,\\
\mathrm{UNRESOLVED} & \text{otherwise.}
\end{cases}
\]
For the credential detector, \code{EXPOSED} records sensitive output without a
known \code{no_log} protection. These labels describe detector evidence; they
are not manual TP/FP judgements.

\section{Rule System and Responsibility Boundaries}
Rules are registered with a stable identifier, category, severity, compatible
artifact classes, detector, message, and evidence fields. The implementation
contains 76 concepts grouped into the canonical result categories. The 38-item
review checklist is represented separately as 15 automated, 18 partial, and 5
manual criteria, because a concept count and checklist coverage answer
different questions.

\begin{table}[H]
\centering\small
\caption{Responsibility of each analysis component.}
\begin{tabularx}{\textwidth}{p{3cm}X X}
\toprule
Component & Delegated responsibility & Deliberately retained elsewhere \\
\midrule
AnsibleGuard & Cross-file variables, execution order, repository structure,
documentation/testing presence, semantic source evidence, canonicalisation &
Does not reproduce every ansible-lint implementation \\
ansible-lint & Established Ansible syntax, FQCN, task/action and formatting
rules under its own configuration and project root & Cross-provider semantic
states and repository-wide evidence \\
GitLeaks & Secret-pattern evidence in content and, when enabled, history &
Ansible execution semantics and final secret classification \\
\bottomrule
\end{tabularx}
\end{table}

Overlap is preserved as provider evidence and merged only when canonical keys
show the same source occurrence and concept. It is not discarded silently.
Provider absence, timeout, non-zero exit, parse failure, or partial output is a
diagnostic with environment context—not an ordinary smell in the analysed
source.

\section{ansible-lint Project Roots}
ansible-lint is invoked at the smallest meaningful project roots. Candidate
roots are normalised and collapsed before execution. If an unconfigured nested
root is covered by an ancestor, it is not linted again; otherwise findings
would be duplicated. Independently configured roots remain separate because
their configuration can change exclusions, profiles, and rule behaviour.
Tests cover both overlapping and independent-root cases.

\section{Canonical Occurrences and Comparisons}
A provider finding is normalised into a versioned occurrence containing source
path, region, concept/rule identity, category, severity, status, message,
fingerprints, and provider evidence. Exact native fingerprints have priority.
Semantic matching requires file identity plus sufficiently specific rule and
location evidence. A generic location-only result is never matched across
files as a rename, avoiding false ``moved'' classifications.

The comparison output partitions occurrences into introduced, fixed,
unchanged, and moved/uncertain sets. Versioned envelopes must pass the shared
v3 validator. The PR comparison utility still accepts the explicitly detected
legacy list form; it does not treat an invalid v3 envelope as legacy data.

\section{Output-v3 Contract}
The canonical JSON Schema is packaged inside the Python distribution and
copied into each scan output. Consequently, an installed wheel validates its
results outside the source checkout. Runtime checks align with the schema for
critical fields: non-empty identifiers and paths, positive coordinates,
bounded regions, recognised severity/status values, typed fingerprints,
provider names and states, and structured provider evidence.

\begin{table}[H]
\centering\small
\caption{Principal scan-output artifacts.}
\begin{tabularx}{\textwidth}{p{4cm}X}
\toprule
Artifact & Purpose \\
\midrule
\code{summary/} & Repository and all-repository aggregates, provider health,
status/severity/category/rule totals, selected files and LOC \\
\code{raw/} & Provider-native findings retained for audit and debugging \\
\code{canonical/} & Deduplicated occurrences with merged provider evidence \\
\code{metadata/} & Commit SHA, configuration, versions, timings, diagnostics,
manifest and reproducibility data \\
\code{files/} & Per-source-file reports where requested \\
\code{diff/} & Introduced, fixed, unchanged and moved/uncertain comparisons \\
\code{output-v3.schema.json} & Exact contract used to validate the output \\
\bottomrule
\end{tabularx}
\end{table}

CSV summaries are projections of the same contract. Repository identity and
SHA establish the population; selected file and LOC columns establish the
denominator; raw and canonical totals distinguish provider output from merged
occurrences; provider columns expose contribution and health; status, severity,
category, and concept columns give orthogonal distributions; diagnostic fields
record incomplete execution. Acronyms are expanded in the report's front
matter, while machine column names remain stable for scripts.

\section{Interfaces and Workflows}
The same engine supports a repository path, a local working-tree change, a
commit range, or a pull request. A full scan discovers all selected artifacts.
Diff modes analyse the relevant revisions and compare canonical occurrences,
so a changed line is not automatically assumed to introduce every finding in
its file. Named profiles select coherent rule/provider sets for full analysis,
security review, correctness, maintainability, CI gating, structural analysis,
or git-focused review. JSON is the authoritative interchange representation;
CSV supports aggregation, SARIF supports code-scanning interfaces, and text or
Markdown supports human review.

\section{Verification Strategy}
Verification is layered: detector unit tests; parsing, scope and graph tests;
provider-root and failure-state tests; schema/serializer contract tests;
comparison regressions; package-install tests; focused remote-equivalence
checks; then the frozen repository study. Documentation staleness,
\code{compileall}, and whitespace checks guard release quality. This verifies
implementation consistency and reproducibility; it is not substituted for a
future independent human accuracy study.
